\documentclass[10pt,letterpaper]{article}

% =========================================================
% Academic CV - faculty search version
% Tailored for geotechnical engineering faculty applications
% =========================================================

\usepackage[left=0.72in,right=0.72in,top=0.62in,bottom=0.58in]{geometry}
\usepackage[T1]{fontenc}
\usepackage[utf8]{inputenc}
\usepackage{newtxtext}
\usepackage{microtype}
\usepackage{xcolor}
\usepackage{tabularx}
\usepackage{array}
\usepackage{enumitem}
\usepackage{titlesec}
\usepackage{hyperref}
\usepackage{fancyhdr}
\usepackage{lastpage}
\usepackage{needspace}
\usepackage{ragged2e}
\usepackage{comment}

\usepackage{etoolbox}


% ---------- Color palette ----------
\definecolor{navy}{HTML}{1F4E68}
\definecolor{charcoal}{HTML}{2E2E2E}
\definecolor{midgray}{HTML}{606060}
\definecolor{rulegray}{HTML}{B7C4CC}

% ---------- Hyperlinks ----------
\hypersetup{
    colorlinks=true,
    urlcolor=navy,
    linkcolor=navy,
    pdfauthor={Srinivas Vivek Bokkisa},
    pdftitle={Academic Curriculum Vitae - Srinivas Vivek Bokkisa}
}

% ---------- Page style ----------
\pagestyle{fancy}
\fancyhf{}
\renewcommand{\headrulewidth}{0pt}
\renewcommand{\footrulewidth}{0pt}
\fancyfoot[L]{\fontsize{8.5}{10}\selectfont\textcolor{midgray}{Srinivas Vivek Bokkisa - Academic CV}}
\fancyfoot[R]{\fontsize{8.5}{10}\selectfont\textcolor{midgray}{Page \thepage\ of \pageref*{LastPage}}}
\setlength{\footskip}{18pt}

% ---------- General typography ----------
\setlength{\parindent}{0pt}
\setlength{\parskip}{0pt}
\setlength{\tabcolsep}{0pt}
\linespread{0.99}\selectfont
\raggedbottom
\emergencystretch=1.5em
\color{charcoal}

\setlist[itemize]{leftmargin=1.15em,itemsep=0.5pt,topsep=1pt,parsep=0pt,partopsep=0pt}
\setlist[enumerate]{leftmargin=1.55em,itemsep=2pt,topsep=1pt,parsep=0pt,partopsep=0pt}


% ---------- Compact section style ----------
\titleformat{\section}
  {\needspace{2\baselineskip}\large\bfseries\color{navy}}
  {}{0pt}{}
  [\vspace{-0.35em}\textcolor{rulegray}{\rule{\textwidth}{0.55pt}}]

\titlespacing*{\section}{0pt}{3pt}{0pt}

% ---------- Compact subsection style ----------
\titleformat{\subsection}
  {\normalsize\bfseries\color{charcoal}}
  {}{1pt}{}

\titlespacing*{\subsection}{0pt}{4pt}{1pt}

% ---------- Reusable commands ----------
\newcommand{\entry}[3]{%
  \begin{tabularx}{\textwidth}{@{}X r@{}}
    \textbf{#1} & \textit{#3}\\[-1pt]
    \textcolor{midgray}{#2} &
  \end{tabularx}\vspace{-1pt}
}


% Degree entry
\newcommand{\degreeentry}[4]{%
    \noindent
    \textbf{#1} \hfill #3\\
    #2%
    \ifstrempty{#4}{}{%
        \enspace|\enspace #4%
    }%
    \par
}

\newcommand{\researchpillar}[2]{%
  \textbf{\textcolor{navy}{#1}}\par
  \vspace{1pt}{\small #2}\par
}


% Numbered publication with an internal hyperlink target
\newcommand{\pub}[2]{%
    \item \hypertarget{pub:#1}{}#2
}

% Clickable citation to a publication
\newcommand{\pubref}[1]{%
    \hyperlink{pub:#1}{\textcolor{navy}{[#1]}}%
}


\begin{document}

% =========================================================
% HEADER
% =========================================================
{\fontsize{20}{22}\selectfont\bfseries\color{charcoal} Srinivas Vivek Bokkisa}\par
\vspace{3pt}
{\large\bfseries\color{navy} Postdoctoral Fellow}\par
\vspace{2pt}
{\small School of Civil and Environmental Engineering, Georgia Institute of Technology | Atlanta, GA, USA}\par
\vspace{5pt}
{\small
\href{mailto:sbokkisa3@gatech.edu}{sbokkisa3@gatech.edu}
\;|\;
\href{https://scholar.google.com/citations?user=DLtSDLgAAAAJ\&hl=en}{Google Scholar}
%\;|\;
%\href{https://orcid.org/0000-0002-0801-3637}{ORCID}
%\;|\;
%\href{https://www.linkedin.com/in/srinivasvivek/}{LinkedIn} 
\;|\;
\href{https://srinivas12viv.github.io/}{Personal Webpage}}


% =========================================================
\section{Research Profile}
% =========================================================
Computational geomechanics researcher advancing
constitutive modeling, numerical methods, and
multiscale geomaterial mechanics to understand and
predict the behavior of geotechnical systems.
Research interests encompass fabric anisotropy,
soil liquefaction and instability, mechanics-informed
site characterization and inverse analysis,
large-deformation simulations, and performance-based
assessment of geotechnical infrastructure, with
emerging directions in scientific machine learning
and AI-informed computational modeling.
\par


% =========================================================
\section{Academic Appointments \& Professional Experience}
% =========================================================

\entry{Postdoctoral Fellow}
{Georgia Institute of Technology,
School of Civil and Environmental Engineering}
{May 2026--Present}
\begin{itemize}
\item Advancing mechanics-based frameworks under a variety of numerical approaches (FEM, MPM, PFEM) for the characterization of complex geomaterials and the assessment of the performance of geosystems (focusing on CPT and slope systems). 

\item Exploring scientific machine learning
and AI-informed approaches for constitutive
modeling and surrogate numerical simulation,
emphasizing the integration of physical principles
to improve predictive generalization.
\end{itemize}

\vspace{2pt}
\entry{Geotechnical Intern -- Mining}{AECOM, Denver, CO, USA}{Jun. 2024--Aug. 2024}
\begin{itemize}
    \item Conducted advanced numerical modeling of tailings storage facilities and developed a database integrating rapid vane shear and co-located CPT data to investigate residual strength in mine tailings.
\end{itemize}



\section{Education}

\degreeentry{Ph.D. in Civil Engineering}
{Georgia Institute of Technology}{2026}
{Advisor: Dr. Jorge Macedo}

\degreeentry{M.S. in Computational Science and Engineering}
{Georgia Institute of Technology}{2023}{}

\degreeentry{M.Phil. in Civil Engineering
(\textit{Concentration in Scientific Computation})}
{Hong Kong University of Science and Technology}{2019}
{Advisor: Dr. Gang Wang}

\degreeentry{B.Tech. in Civil Engineering}
{Indian Institute of Technology Guwahati}{2017}{}

% =========================================================
\section{Awards \& Honors}
% =========================================================
\begin{itemize}[leftmargin=*, itemsep=2pt, topsep=2pt]
    \item \textbf{George F. Sowers Distinguished Graduate
    Student Award}, Georgia Tech, 2026.
    Recognizes outstanding doctoral achievement in
    scholarship, research, teaching, and service.

    \item \textbf{Outstanding Geomaterials Research Award},
    Georgia Tech, 2025.
    Recognizes research novelty and original
    contributions to geotechnical engineering.

    \item \textbf{Selected Participant, ALERT Olek
    Zienkiewicz School on Constitutive Modeling of
    Geomaterials}, Prague, Czechia, 2025.

    \item \textbf{Paul W. Mayne Award}, Georgia Tech, 2024.
    Recognizes outstanding academic performance,
    research, and service during early doctoral study.

    \item \textbf{Best Poster Award}, Annual Geotechnical
    Student Poster Symposium, Geo-Institute
    \& Georgia Tech, 2024.

    \item \textbf{Selected Participant, APRU-IRIDeS
    Multi-Hazards Summer School},
    Tohoku University, Sendai, Japan, 2018.

    \item \textbf{Postgraduate Scholarship},
    Hong Kong University of Science and Technology,
    2017--2019.
\end{itemize}

% =========================================================
\section{Publications \& Manuscripts}
% =========================================================
% =========================================================
\subsection{Peer-Reviewed Journal Articles}
% =========================================================

\begin{enumerate}[label={[\arabic*]}, leftmargin=*, itemsep=0pt]

\pub{1}{\textbf{Srinivas Vivek, B.}, and Macedo, J. (2026). ``Fabric Effects on Cone Penetration Testing: A Numerical Investigation.'' \textit{Computers and Geotechnics}. Accepted. \href{https://papers.ssrn.com/sol3/papers.cfm?abstract_id=6864496}{[preprint]}}

\pub{2}{\textbf{Srinivas Vivek, B.}, Macedo, J., and Arduino, P. (2026). ``On the integration of an ACST-based bounding surface model.'' \textit{Computers and Geotechnics}. \href{https://doi.org/10.1016/j.compgeo.2026.107911}{[DOI]}}

\pub{3}{\textbf{Srinivas Vivek, B.}, and Macedo, J. (2025). ``Challenges in NorSand to model CSD stress paths and proposed modifications.'' \textit{Canadian Geotechnical Journal}. \href{https://cdnsciencepub.com/doi/pdf/10.1139/cgj-2024-0546}{[DOI]}}

\pub{4}{\textbf{Srinivas Vivek, B.}, and Macedo, J. (2024). ``Influence of anisotropic consolidation on the instability of loose granular soils under undrained and drained loading.'' \textit{Canadian Geotechnical Journal}. \href{https://cdnsciencepub.com/doi/abs/10.1139/cgj-2023-0683}{[DOI]}}

\pub{5}{\textbf{Srinivas Vivek, B.}, Macedo, J., and Petalas, A. L. (2023). ``Fabric anisotropy effects on static liquefaction under constant shear drained loading.'' \textit{Computers and Geotechnics}. \href{https://www.sciencedirect.com/science/article/pii/S0266352X23004810}{[DOI]}}

\pub{6}{\textbf{Srinivas Vivek, B.}, Macedo, J., Petalas, A. L., and Arson, C. (2022). ``Assessing flow liquefaction triggering considering fabric anisotropy effects under the ACST framework.'' \textit{Computers and Geotechnics}. \href{https://www.sciencedirect.com/science/article/pii/S0266352X22001537}{[DOI]}}

\pub{7}{Kumar, A., and \textbf{Srinivas, B. V.} (2017). ``Easy to use empirical correlations for liquefaction and no liquefaction conditions.'' \textit{Geotechnical and Geological Engineering}, 35(4), 1383--1407. \href{https://link.springer.com/content/pdf/10.1007\%2Fs10706-017-0183-2.pdf}{[DOI]}}

\end{enumerate}


% =========================================================
\subsection{Manuscripts Ready for Submission}
% =========================================================
\begin{enumerate}[resume, label={[\arabic*]}, leftmargin=*, itemsep=0pt]

\pub{8}{\textbf{Srinivas Vivek, B.}, and Macedo, J. ``Directional Second-Order Work and Instability under Triaxial Stress Relief.'' Manuscript ready for submission; under internal review.}

\pub{9}{\textbf{Srinivas Vivek, B.}, and Macedo, J. ``Fabric Effects on Slope Instability and Deformation:
A Numerical Investigation.'' Manuscript ready for submission; under internal review.}

\end{enumerate}

% =========================================================
\section{Conference Papers \& Presentations}
% =========================================================
\begin{enumerate}[resume, label={[\arabic*]}, leftmargin=*, itemsep=0pt]

\pub{10}{\textbf{Srinivas Vivek, B.}, and Macedo, J. (2025). ``Challenges in NorSand for modeling stress-relief paths in tailings storage facilities.'' \textit{Tailings and Mine Waste Conference 2025}.}

\pub{11}{\textbf{Srinivas Vivek, B.}, and Macedo, J. (2024). ``Fabric effects on the onset of static liquefaction under drained conditions: Implications for mine tailings.'' \textit{Tailings and Mine Waste Conference 2024}.}

\pub{12}{\textbf{Srinivas Vivek, B.}, Macedo, J., and Petalas, A. L. (2023). ``On the effects of fabric on the instability onset under constant shear drained loading.'' \textit{EMI Conference 2023}.}

\pub{13}{\textbf{Srinivas Vivek, B.}, Macedo, J., Petalas, A. L., and Arson, C. (2023). ``Effects of initial consolidation on the triggering of static liquefaction considering fabric effects.'' \textit{Geo-Congress 2023}. \href{https://ascelibrary.org/doi/epdf/10.1061/9780784484708.008}{[PDF]}}

\pub{14}{\textbf{Srinivas Vivek, B.}, Macedo, J., Petalas, A. L., and Arson, C. (2022). ``Effects of soil fabric on the triggering of static liquefaction: Implications for mine tailings.'' \textit{Tailings and Mine Waste Conference 2022}.}

\pub{15}{\textbf{Srinivas Vivek, B.}, Macedo, J., Petalas, A. L., and Arson, C. (2022). ``On the effects of fabric on the instability surface of granular materials.'' \textit{EMI Conference 2022}.}

\pub{16}{\textbf{Srinivas Vivek, B.}, Wang, G., Duruo, H., and Feng, J. (2019). ``Fabric evolution in post-liquefaction and re-liquefaction of granular soils using 3D discrete element modeling.'' \textit{Proceedings of the 7th International Conference on Earthquake Geotechnical Engineering}, Rome, Italy, vol. 4, pp. 1461--1468. CRC Press.}

\pub{17}{\textbf{Srinivas Vivek, B.}, and Wang, G. (2019). ``DEM modeling of Toyoura sand considering real particle shapes.'' \textit{23rd Annual Conference of HKSTAM and 15th Jiangsu-Hong Kong Forum on Mechanics and its Application}, Hong Kong, p. 41.}

\pub{18}{\textbf{Srinivas Vivek, B.}, and Mohanty, S. (2016). ``1D ground response analysis of Bhubaneswar soil in India.'' \textit{6th International Conference on Recent Advances in Geotechnical Earthquake Engineering and Soil Dynamics}, New Delhi, India.}

\end{enumerate}

% =========================================================
\section{Theses \& Dissertations}
% =========================================================
\begin{enumerate}[resume, label={[\arabic*]}, leftmargin=*, itemsep=0pt]

\pub{19}{\textbf{Srinivas Vivek, B.} (2026).
``Multiscale Assessment of Fabric in Granular Soils.''
Ph.D. dissertation, Georgia Institute of Technology.}

\pub{20}{\textbf{Srinivas Vivek, B.} (2019).
``Fabric Evolution in 3D Discrete Element modeling of
Soil Liquefaction.''
M.Phil. dissertation, Hong Kong University of Science
and Technology. \href{https://doi.org/10.14711/thesis-991012752454703412}{[DOI]}}

\end{enumerate}


% =========================================================
\section{Research Software \& Computational Tools}
% =========================================================

\entry{Microstructure \& Fabric Explorer ---
Granular Fabric Analysis Toolkit}
{Scientific Toolkit Development |
\href{https://srinivas12viv.github.io/microstructure-fabric-explorer/}
{Project Webpage}}
{2019--Present}
\begin{itemize}
    \item Developed computational tools for granular
    microstructure characterization using DEM simulations
    and X-ray micro-CT images, integrating contact detection,
    pore-space segmentation, fabric tensor analysis,
    force-network characterization, and Minkowski-based
    particle--void descriptors with interactive 2D/3D
    visualizations to investigate microstructural
    mechanisms underlying geomaterial behavior.
\end{itemize}

\vspace{2pt}
\entry{SoilWorkspace --- Integrated Soil Data Analysis,
Calibration, and Simulation Platform}
{Independent Software Development | Lead Developer |
\href{https://srinivas12viv.github.io/SoilWorkspace/}
{Project Glimpse}}
{May 2026--Present}
\begin{itemize}
    \item Developing an extensible platform integrating
    geotechnical laboratory data analysis, constitutive
    model calibration, and virtual element testing,
    with reproducible workflows, support for user-defined
    models, and interoperability with simulation software
    to bridge research and engineering practice.

    \item Planned extensions include automatic-differentiation-based
    frameworks for constitutive model implementation
    and verification, simplifying model development
    and enabling systematic evaluation of formulations
    and numerical integration schemes.
\end{itemize}

\vspace{2pt}
\entry{NorSand Interactive Visualizations ---
Constitutive Modeling Educational Resource}
{Computational Education |
\href{https://srinivas12viv.github.io/NorSand/}
{Interactive Resource}}
{2022}
\begin{itemize}
    \item Developed interactive educational visualizations
    of NorSand constitutive mechanics, illustrating image
    state, dilatancy, yield-surface evolution, and hardening
    through interactive plots to strengthen
    students' physical intuition and understanding of
    critical state soil mechanics; used in instruction
    at Georgia Tech, Colorado School of Mines, and
    Virginia Tech.
\end{itemize}



% =========================================================
\section{Scientific Computing \& Technical Skills}
% =========================================================

\textbf{Programming:}
Python, Fortran, C++, MATLAB.
\vspace{2pt}

\textbf{Numerical Methods \& Software:}
DEM (LIGGGHTS, YADE);
FEM (PLAXIS, Abaqus, OpenSees);
MPM (GeoTaichi, Anura3D);
PFEM (Kratos); constitutive model implementation
(UMAT/UDSM), numerical integration and verification.
\vspace{2pt}

\textbf{High-Performance Computing \& Cyberinfrastructure:}
TACC, DesignSafe, Slurm.
\vspace{2pt}

\textbf{Scientific Computing \& Visualization:}
ParaView, ImageJ, SPAM;
X-ray micro-CT image analysis and fabric characterization.

% =========================================================
\section{Teaching Experience}
% =========================================================
\entry{Graduate Teaching Assistant}{Georgia Institute of Technology}{Jan. 2021--May 2025}
\begin{itemize}
    \item \textbf{Courses:} Introduction to Geotechnical Engineering; Data Analytics in Civil and Environmental Engineering.
    \item Developed and delivered tutorials on machine learning, AI, and data analytics using civil engineering datasets.
    \item Led two laboratory sections, conducted geotechnical laboratory demonstrations, held office hours, and evaluated student assignments.
\end{itemize}

\vspace{2pt}
\entry{Graduate Teaching Assistant}{Hong Kong University of Science and Technology}{Feb. 2018--May 2019}
\begin{itemize}
    \item Supported instruction in \textit{Rock Mechanics} and \textit{Construction Materials} through tutorial teaching, concrete laboratory demonstrations, assignment evaluation, and grading.
\end{itemize}

\vspace{5pt}
\textbf{Teaching Interests:}
Computational Geomechanics; Soil Mechanics;
Critical State Soil Mechanics; Constitutive Modeling;
Numerical Methods in Geotechnical Engineering;
Geotechnical Earthquake Engineering;
Site Characterization; Scientific Machine Learning
for Geotechnical Engineering.


% =========================================================
\section{Student Mentoring}
% =========================================================

\entry{Graduate Research Mentor, Dr. Jorge Macedo's Research Group}
{Georgia Institute of Technology}{Aug. 2023--Present}
\begin{itemize}
    \item Mentored graduate students in computational
    geomechanics, including a master's student who
    transitioned to doctoral research at Caltech
    and a current Ph.D. student. Provided guidance
    in constitutive theory, numerical implementation,
    and interpretation of simulation results,
    emphasizing independent problem-solving,
    critical thinking, and reproducible research
    practices.
\end{itemize}


\entry{VIP Mentor, DeepMeLT and Extreme Events Engineering Groups}{Georgia Institute of Technology}{Sep. 2021--May 2025}
\begin{itemize}
    \item Mentored 4--6 undergraduate researchers per semester
    in computational geomechanics
    and high-performance computing using TACC resources,
    fostering independent problem-solving, computational
    proficiency, and research skills.
\end{itemize}

\vspace{2pt}
\entry{Student Academic Mentor, Geosystems Engineering}{Georgia Institute of Technology}{Jan. 2024--Dec. 2024}
\begin{itemize}
    \item Mentored a first-year graduate student through monthly meetings focused on academic progress, professional development, and graduate-community integration.
\end{itemize}

\vspace{2pt}
\entry{Student Academic Mentor}{Indian Institute of Technology Guwahati}{Aug. 2015--Apr. 2016}

% =========================================================
\section{Invited Lectures \& Seminars}
% =========================================================
\begin{itemize}
    \item \textbf{Invited expert lecture}, ``Toward Performance-Based Assessment in Tailings Engineering,'' Indian Geotechnical Society Dhanbad chapter, IIT (ISM) Dhanbad, India, 2026.

    \item \textbf{Invited speaker}, ``Fabric Effects in Static Liquefaction,'' Barr Engineering Co., 2023.

    \item \textbf{Invited lectures}, ``Visualizing NorSand,'' Critical State Soil Mechanics courses, Colorado School of Mines and Georgia Tech.
\end{itemize}
% =========================================================
\section{Academic \& Professional Service}
% =========================================================
\textbf{Peer Review Service}\par
\vspace{2pt}
Reviewer for \textit{JGGE} (ASCE), \textit{IJG} (ASCE), and \textit{Open Geomechanics}.
\vspace{2pt}

\entry{Vice President, Geotechnical Society}{Georgia Institute of Technology}{Aug. 2023--Aug. 2024}
\begin{itemize}
    \item Managed organizational finances and helped organize academic and community events for the Geosystems community.
\end{itemize}
\vspace{2pt}

\entry{General Secretary of Academics, Association of Civil Engineers}{Indian Institute of Technology Guwahati}{Oct. 2015--Oct. 2016}
\end{document}